% =============================================================
% TER -- Pacman Labyrinth: Mini Report on AI Management
% Self-contained LaTeX file. Compile with:  pdflatex report.tex
% =============================================================
\documentclass[11pt,a4paper]{article}

% ---------- Encoding & language ----------
\usepackage[T1]{fontenc}
\usepackage[utf8]{inputenc}
\usepackage[english]{babel}
\usepackage{lmodern}
\usepackage{microtype}

% ---------- Layout ----------
\usepackage[a4paper,margin=2.2cm]{geometry}
\usepackage{titlesec}
\usepackage{enumitem}
\usepackage{parskip}
\usepackage{multicol}

% ---------- Visual / utility ----------
\usepackage{xcolor}
\usepackage{graphicx}
\usepackage{tcolorbox}
\tcbuselibrary{skins,breakable}
\usepackage{booktabs}
\usepackage{array}
\usepackage{tabularx}
\usepackage{listings}
\usepackage{amsmath}
\usepackage{amssymb}
\usepackage{hyperref}
\usepackage{fancyhdr}

% ---------- Colours (Pac-Man inspired) ----------
\definecolor{pmblue}{RGB}{33,33,222}
\definecolor{pmyellow}{RGB}{255,204,0}
\definecolor{pmpink}{RGB}{255,184,222}
\definecolor{pmcyan}{RGB}{0,200,222}
\definecolor{pmorange}{RGB}{255,140,0}
\definecolor{pmred}{RGB}{222,40,40}
\definecolor{pmgreen}{RGB}{40,180,90}
\definecolor{darkbg}{RGB}{20,20,40}
\definecolor{softgrey}{RGB}{245,245,250}
\definecolor{rulegrey}{RGB}{200,200,210}

% ---------- Hyperlinks ----------
\hypersetup{
    colorlinks=true,
    linkcolor=pmblue,
    urlcolor=pmblue,
    citecolor=pmblue
}

% ---------- Section style ----------
\titleformat{\section}
  {\color{pmblue}\Large\bfseries\sffamily}
  {\thesection}{0.6em}{}
  [\vspace{-0.3em}{\color{pmblue}\rule{\linewidth}{0.6pt}}]
\titleformat{\subsection}
  {\color{darkbg}\large\bfseries\sffamily}
  {\thesubsection}{0.6em}{}
\titleformat{\subsubsection}
  {\color{darkbg}\normalsize\bfseries\sffamily}
  {\thesubsubsection}{0.5em}{}
\titlespacing*{\section}{0pt}{1.0em}{0.4em}
\titlespacing*{\subsection}{0pt}{0.7em}{0.2em}
\titlespacing*{\subsubsection}{0pt}{0.5em}{0.15em}

% ---------- Header / footer ----------
\pagestyle{fancy}
\fancyhf{}
\renewcommand{\headrulewidth}{0.4pt}
\renewcommand{\headrule}{\color{rulegrey}\hrule height 0.4pt}
\fancyhead[L]{\sffamily\small\color{pmblue}TER -- Pacman Labyrinth}
\fancyhead[R]{\sffamily\small\color{pmblue}AI Management Report}
\fancyfoot[C]{\sffamily\small\thepage}

% ---------- Code listings ----------
\lstdefinestyle{mini}{
  basicstyle=\ttfamily\footnotesize,
  keywordstyle=\color{pmblue}\bfseries,
  commentstyle=\color{gray}\itshape,
  stringstyle=\color{pmorange},
  numbers=left,
  numberstyle=\tiny\color{gray},
  numbersep=8pt,
  backgroundcolor=\color{softgrey},
  frame=single,
  rulecolor=\color{rulegrey},
  framesep=4pt,
  breaklines=true,
  showstringspaces=false,
  tabsize=2,
  captionpos=b,
}
\lstset{style=mini}

% ---------- Custom boxes ----------
\newtcolorbox{infobox}[1][]{
  enhanced, breakable,
  colback=pmblue!5, colframe=pmblue,
  arc=2pt, boxrule=0.6pt,
  left=8pt,right=8pt,top=4pt,bottom=4pt,
  fonttitle=\bfseries\sffamily\color{white},
  title=#1
}
\newtcolorbox{ghostbox}[2][]{
  enhanced, breakable,
  colback=#2!8, colframe=#2,
  arc=2pt, boxrule=0.6pt,
  left=6pt,right=6pt,top=3pt,bottom=3pt,
  fonttitle=\bfseries\sffamily\color{white},
  title=#1
}

% =============================================================
\begin{document}

% ---------- Cover ----------
\begin{center}
{\color{pmblue}\rule{\linewidth}{1.2pt}}\\[0.4em]
{\sffamily\Huge\bfseries\color{darkbg} Pacman Labyrinth}\\[0.3em]
{\sffamily\Large\color{pmblue} A Mini Report on AI Management}\\[0.4em]
{\color{pmblue}\rule{\linewidth}{0.6pt}}\\[1.2em]

\begin{tabular}{rl}
\textbf{Project} & TER -- Master 1 IA \\
\textbf{Group}   & TER\_S2\_G \\
\textbf{Members} & Daniel Abou Orm \\
                 & Joe Kabbat \\
                 & Kacem Abdelbaki \\
\textbf{Date}    & May 2026 \\
\end{tabular}

\vspace{0.8em}
\begin{infobox}[Scope]
\textit{This document explains how artificial intelligence is handled in
our Pacman--Labyrinth project: ghost behaviour, the autonomous Pac-Man
bot, the difficulty system, the supporting layers (maze generation,
rendering, UI, networking, replay), and a final reflection on our use of
AI \emph{tools} (Copilot, ChatGPT) during development.}
\end{infobox}
\end{center}

% ---------- 1 ----------
\section{Project Overview}

The project is a complete Pacman--Labyrinth platform split into two
cooperating components communicating exclusively through a REST API.

\begin{itemize}[leftmargin=1.3em,itemsep=2pt]
  \item \textbf{Backend (\texttt{labyrinth\_gen}, Python / Flask).}
        Generates symmetric mazes using four algorithms
        (\textsc{Dfs}, \textsc{Kruskal}, \textsc{Prim}, \textsc{Tetris}),
        persists them in MongoDB, and exposes a REST API documented with
        Swagger. Containerised with Docker and deployed on Render.
  \item \textbf{Frontend (\texttt{labyrinth\_ui}, Java / LibGDX).}
        Fetches a maze from the API, renders it in a Pac-Man style, and
        runs the gameplay: human or bot Pac-Man, four classic ghosts,
        dots, power pellets, teleporters, and a replay system.
\end{itemize}

\begin{infobox}[Why this matters for the AI report]
The pipeline forces every AI decision to operate on a \emph{procedurally
generated} grid that is never the same twice. Ghost targeting and the
Pac-Man bot must therefore reason on the fly, with no hand-tuned map.
\end{infobox}

% ---------- 2 ----------
\section{Maze Generation: Foundation of the AI World}

Before any agent can think, the world must exist. The four generators
share a common pipeline (\texttt{common.py}):

\begin{center}\small
\textit{init grid} $\rightarrow$ \textit{carve} $\rightarrow$
\textit{add cycles} $\rightarrow$ \textit{mirror for symmetry} $\rightarrow$
\textit{insert middle column} $\rightarrow$ \textit{place teleporters}
$\rightarrow$ \textit{seal borders}.
\end{center}

The output is a serialisable \texttt{Cell} grid. Each cell carries four
boolean walls (\texttt{[Up, Right, Down, Left]}), an \texttt{isGhost}
flag for the ghost room, and an \texttt{isTP} flag (with a target
\texttt{TPCell}) for teleport pairs. This compact representation is what
both the Pac-Man bot and the ghosts will read at every tick.

\begin{table}[h]
\centering\small
\renewcommand{\arraystretch}{1.15}
\begin{tabularx}{\linewidth}{l X}
\toprule
\textbf{Algorithm} & \textbf{Behaviour and impact on AI} \\
\midrule
\textsc{Dfs}     & Recursive backtracking. Long winding corridors, few
                   junctions. Ghosts feel predictable and the bot's path is
                   mostly linear. \\
\textsc{Kruskal} & Spanning tree via Union-Find on edges. Uniform branching
                   with many short choices, harder for a one-step greedy
                   bot. \\
\textsc{Prim}    & Frontier-based growth. Compact ``bushy'' areas;
                   ghosts converge faster on Pac-Man. \\
\textsc{Tetris}  & Active-cell growth with branching bias ($p\!=\!0.85$).
                   Produces large rooms and bottlenecks; the most playable
                   and balanced for both humans and bots. \\
\bottomrule
\end{tabularx}
\end{table}

\subsection{Symmetry, cycles and teleporters}

Three post-processing passes shape the AI experience as much as the
underlying algorithm:

\begin{itemize}[leftmargin=1.3em,itemsep=1pt]
  \item \textbf{Mirroring} (\texttt{mirror\_map}) doubles the carved
        half-grid into a left/right symmetric maze. Symmetry matters: it
        lets us place the ghost room and the four scatter corners
        coherently regardless of the random carving.
  \item \textbf{Cycles} (\texttt{add\_cycles(p)}) randomly knock down
        walls so the graph is no longer a tree. Without cycles, ghosts
        would always have a unique path to Pac-Man (chase becomes
        trivial) and the bot would have no escape route.
  \item \textbf{Teleporters} (\texttt{add\_TP\_Points}) add two pairs of
        portal cells that introduce \emph{non-grid} edges in the maze
        graph. Both ghosts and the bot must explicitly recognise these
        edges; we will see in \S4 that they were the source of an
        oscillation bug in the bot.
\end{itemize}

A small \texttt{maze\_benchmarks} routine counts cycles and hallway
lengths so we could tune \texttt{p} empirically and confirm that
\textsc{Tetris} produces the highest playability score.

% ---------- 3 ----------
\section{Pacman AI: Ghost Behaviour}

Ghosts are the central AI artefact of the project. We follow the
\textbf{classic Pac-Man targeting model} \cite{pacmandossier} and extend
it with an instrumented \emph{A*} mode for academic comparison.

\subsection{Architecture and shared mechanics}

All ghosts inherit from an abstract \texttt{Ghost} class which
implements the shared movement contract. The contract is intentionally
rigid so each subclass only has to specify \emph{where to aim}, not
\emph{how to move}.

\begin{itemize}[leftmargin=1.3em,itemsep=1pt]
  \item \textbf{Four modes:} \texttt{CHASE} (hunt Pac-Man),
        \texttt{SCATTER} (retreat to a corner),
        \texttt{FRIGHTENED} (random walk after Pac-Man eats a power
        pellet) and \texttt{EATEN} (return to the ghost room as eyes).
  \item \textbf{Discrete grid + smooth visuals.} Logical moves happen on
        a fixed grid every \texttt{moveInterval} seconds (default
        $0.22\,$s). Visual position is driven by a per-frame
        \texttt{MathUtils.lerp} between the current and the target tile
        so the AI's discrete tick is hidden.
  \item \textbf{Forbidden reverse.} A ghost never turns
        180\textdegree\ unless it is the only legal move. This single
        rule, lifted from the original arcade, is what creates most of
        the emergent ``personality''.
  \item \textbf{Tie-break priority} \(\textsc{Up} > \textsc{Left} >
        \textsc{Down} > \textsc{Right}\), preserving the famous arcade
        quirk that lets advanced players read ghost behaviour.
  \item \textbf{Two difficulty modes:} \texttt{CLASSIC} (greedy
        Euclidean step) and \texttt{A\_STAR} (true shortest-path
        planning). Both share the rest of the contract.
\end{itemize}

The mode timeline is driven from the screen by
\texttt{MazeScreen.MODE\_DURATIONS = \{7, 20, 7, 20, 5, 20, 5, $\infty$\}}
seconds, alternating SCATTER and CHASE -- the original 1980 Namco
schedule. \texttt{FRIGHTENED} is a temporary overlay (6 seconds) on top
of that timeline, triggered when the player eats a power pellet.

\subsubsection{Coordinate conventions}

A subtle but recurrent source of bugs was the dual coordinate convention
between cells and movement: cell walls are indexed
\texttt{[Up, Right, Down, Left]} but agent directions are indexed
\(\texttt{Right}{=}0,\,\texttt{Up}{=}1,\,\texttt{Left}{=}2,\,
\texttt{Down}{=}3\). The \texttt{Ghost} class therefore keeps an
explicit mapping table \texttt{DIR\_TO\_WALL[]\,=\,\{1,0,3,2\}} so that
direction-to-wall lookups stay correct even after refactoring.

\subsection{The four personalities}

\begin{ghostbox}[Blinky -- The Chaser]{pmred}
Targets Pac-Man's exact tile in \texttt{CHASE} mode. Scatter corner:
top-right. Simple, relentless; sets the baseline pressure and serves as
a reference position for Inky's vector trick.
\end{ghostbox}

\begin{ghostbox}[Pinky -- The Ambusher]{pmpink}
Aims four tiles \emph{ahead} of Pac-Man in his current direction.
Encourages the player to make sharp turns to break the trap. Scatter
corner: top-left.
\end{ghostbox}

\begin{ghostbox}[Inky -- The Unpredictable]{pmcyan}
Computes a vector from Blinky to a point two tiles ahead of Pac-Man,
then \emph{doubles} it. Inky's target therefore depends on
\emph{another ghost's position}, producing the cooperative pincer
behaviour the original arcade is famous for. Scatter corner:
bottom-right.
\end{ghostbox}

\begin{ghostbox}[Clyde -- The Shy One]{pmorange}
Chases Pac-Man only when he is more than 8 tiles away (Euclidean);
otherwise retreats to his scatter corner. Adds a non-monotonic, almost
timid agent that prevents the group dynamics from collapsing into
``four red Blinkies''. Scatter corner: bottom-left.
\end{ghostbox}

\subsection{The greedy heuristic (\texttt{CLASSIC} mode)}

In \texttt{CLASSIC} mode the ghost performs a one-step lookahead.
Among the (up to three) non-reverse neighbours that are not separated by
a wall, it picks the one minimising the \emph{squared Euclidean
distance} to its target tile, with ties broken by the priority listed
above. The move is computed in $O(1)$, faithful to the arcade, but
locally myopic -- a ghost can be lured into dead-ends if Pac-Man dances
around a wall right behind it.

\subsection{The A* heuristic (\texttt{A\_STAR} mode)}

\texttt{A\_STAR} mode runs a full A* search over the maze graph from
the ghost's current cell to its target tile, using the Manhattan
distance as an admissible heuristic. Only the \emph{first move} of the
optimal path is returned -- A* is recomputed every tick because Pac-Man
moves and Inky's reference (Blinky) shifts every frame.

\begin{lstlisting}[language=Java,caption={A* in \texttt{Ghost.chooseDirectionAStar} -- abridged.}]
PriorityQueue<Node> open = new PriorityQueue<>();
boolean[][] visited = new boolean[maze.height][maze.width];

// Seed the open set with the (up to 3) non-reverse first moves
int reverseDir = (direction + 2) % 4;
for (int d = 0; d < 4; d++) {
    if (d == reverseDir || startCell.hasWall(DIR_TO_WALL[d])) continue;
    int nr = gridRow + DR[d], nc = gridCol + DC[d];
    int h  = Math.abs(nr - targetRow) + Math.abs(nc - targetCol);
    open.add(new Node(nr, nc, /*g*/ 1, /*h*/ h, /*firstDir*/ d));
}
while (!open.isEmpty()) {
    Node cur = open.poll();
    if (cur.r == targetRow && cur.c == targetCol) return cur.dir;
    if (visited[cur.r][cur.c]) continue;
    visited[cur.r][cur.c] = true;
    /* expand 4 neighbours, propagating cur.dir as the first move */
}
\end{lstlisting}

A key implementation detail is that \texttt{firstDir} (the direction of
the very first step) propagates unchanged through the search tree: when
A* finally pops the goal, the action returned is the move that started
the optimal path. This avoids reconstructing the full path.

\subsection{Mode-specific behaviour}

\texttt{FRIGHTENED} ghosts deliberately bypass A* (they are supposed to
be \emph{less} smart, not more) and revert to a randomised classic
move. \texttt{EATEN} ghosts ignore walls altogether and home back to
their spawn via \texttt{moveTowardGhostRoom}, then exit again through
the door cell detected dynamically with \texttt{exitGhostRoom}. We
added this dynamic door detection because the door's column varies with
the maze size and with the mirroring step.

\subsection{Telemetry and measurements}

To compare both modes objectively we instrumented every ghost decision:

\begin{lstlisting}[language=Java,caption={Per-decision telemetry on every ghost.}]
protected long lastComputeTimeNs = 0;
protected int  lastNodesEvaluated = 0;
\end{lstlisting}

The HUD reads these counters and prints the live ghost cost on screen.
Averaged over a 60-second \texttt{RaceScreen} run on a $25\!\times\!25$
Tetris-generated maze, we observed:

\begin{table}[h]
\centering\small
\renewcommand{\arraystretch}{1.15}
\begin{tabular}{lrrr}
\toprule
\textbf{Mode} & \textbf{Avg.\ time / decision} & \textbf{Nodes evaluated} &
\textbf{Catch rate vs.\ greedy bot} \\
\midrule
\texttt{CLASSIC} & $\approx\!1\,\mu$s   & $\le 3$            & $\sim$45\% \\
\texttt{A\_STAR} & $\approx\!120\,\mu$s & $\sim\!200$--$400$ & $\sim$78\% \\
\bottomrule
\end{tabular}
\end{table}

Both modes stay well below the 16.6\,ms frame budget at 60 FPS, even
with four concurrent ghosts. We therefore expose \texttt{A\_STAR} as a
configurable ``hard mode'' from the menu.

% ---------- 4 ----------
\section{Pacman AI: The Autonomous Bot}

The Pac-Man side also carries a non-trivial AI. \texttt{BotAgent} can
play a full game without human input, and is also what we use to
benchmark the ghost difficulty levels of \S3. It implements the same
\texttt{Agent} interface as \texttt{HumanAgent}, so the rest of the
engine (movement, collision, scoring, replay) is unchanged.

\subsection{Game-state injection}

Each frame, the screen calls \texttt{botAgent.setGameState(dots,
ghosts, playerRow, playerCol)}. The bot does \emph{not} keep references
to live game objects beyond this call; instead it stores a snapshot,
which makes the planner deterministic and easy to test. The same call
also triggers a small piece of bookkeeping: if Pac-Man just landed on a
teleport tile, the bot starts a TP cooldown (see \S4.4).

\subsection{Greedy strategy: BFS to the nearest dot}

The first strategy, \texttt{BotStrategy.GREEDY}, runs a breadth-first
search from Pac-Man's tile and returns the first step of the shortest
path to the closest uneaten dot. BFS is preferred over a naive
``closest by Euclidean distance'' because it accounts for walls -- the
visually closest dot is often unreachable in a few moves.

\begin{lstlisting}[language=Java,caption={Greedy BFS dot-seeking (\texttt{BotAgent.computeGreedyMove}, abridged).}]
Queue<int[]> q = new LinkedList<>();
visited[playerRow][playerCol] = true;
firstStep[playerRow][playerCol] = -1;

for (int d = 0; d < 4; d++) {                          // seed the 4 neighbours
    int nr = playerRow + DIRS[d][0], nc = playerCol + DIRS[d][1];
    if (canMove(playerRow, playerCol, d)
        && !isTpCooldown(nr, nc)) {
        visited[nr][nc] = true;
        firstStep[nr][nc] = d;
        if (dots[nr][nc]) return new int[]{ DIRS[d][0], DIRS[d][1] };
        q.add(new int[]{ nr, nc });
    }
}
while (!q.isEmpty()) { /* expand neighbours + follow teleport edges */ }
\end{lstlisting}

The trick that keeps the implementation small is the
\texttt{firstStep[]} array: every cell stores the direction of the
\emph{first} step taken from Pac-Man to reach it, so when BFS finally
discovers a dot we recover the move in $O(1)$ without rebuilding the
path. Teleport edges are processed in the same loop, which keeps the
state space connected.

\subsection{Ghost-avoidance strategy: Dijkstra over a danger field}

The second strategy, \texttt{BotStrategy.GHOST\_AVOIDANCE}, is the
heart of the bot. Inspired by \emph{influence maps}
\cite{millington}, we precompute a danger field by spreading a penalty
from every active ghost via a bounded BFS:
\[
\text{danger}(r,c) \;=\; \sum_{g\,\in\,\text{active ghosts}}
\max\!\bigl(0,\; R - d_{\text{maze}}(g,(r,c))\bigr),
\quad R = 4.
\]

\texttt{FRIGHTENED} and \texttt{EATEN} ghosts are deliberately excluded
from the field -- they are not threats, and avoiding them would
actually prevent the bot from ``eating ghosts'' for bonus points. The
field is then used as an additive cost on a Dijkstra search whose
target is the nearest dot:
\[
\text{cost}(u\!\to\!v) \;=\; 1 + 3\cdot \text{danger}(v).
\]

The factor $3$ was tuned empirically: lower values let the bot walk
straight into ghosts; higher values made it freeze in corners far from
any dot. The Dijkstra search uses the same \texttt{firstStep[]} trick
as BFS, with the additional rule that the \emph{first} step is locked
in only when leaving Pac-Man's own cell -- a small but easy-to-miss
invariant that we documented inline.

\subsubsection{Fallback}

If Dijkstra cannot find any dot (for instance because the danger field
\emph{surrounds} Pac-Man), the bot falls back to the greedy BFS. We
prefer ``run for a dot, possibly through danger'' over standing still
-- standing still always loses.

\subsection{Subtle detail: teleport handling}

\begin{infobox}[The teleport oscillation bug]
Teleport tiles introduce non-grid edges. We extend both BFS and
Dijkstra to follow them, but during early playtests the bot would
sometimes \textbf{oscillate inside a teleport pair}: it would step
through portal $A\!\to\!B$, immediately notice that the closest dot is
on the side of $A$, walk back $B\!\to\!A$, and so on. The fix is a
\textbf{TP cooldown}: when the bot arrives on a TP exit cell, that cell
is forbidden for the next 5 ticks. The cooldown is consumed inside
\texttt{setGameState}, so it survives across both BFS and Dijkstra
invocations.
\end{infobox}

\subsection{Performance and effective playstyle}

On a $25\!\times\!25$ maze with four ghosts the worst case is the
ghost-avoidance strategy: BFS for danger ($\le 4$ ghosts $\times$
$\le 200$ nodes) plus Dijkstra ($\le 625$ nodes). The total stays
sub-millisecond on a laptop CPU and is run only once per Pac-Man tick
($\approx 130\,$ms), so the bot is essentially free.

Qualitatively, the two strategies behave very differently:

\begin{itemize}[leftmargin=1.3em,itemsep=1pt]
  \item \textbf{Greedy} clears small mazes quickly but is suicidally
        confident -- it walks straight into a ghost if a dot happens to
        be on the other side. We use it as a stress test for the ghosts.
  \item \textbf{Ghost-avoidance} survives much longer and exhibits
        emergent behaviour: it reverses direction at junctions when a
        ghost is entering, hugs walls to avoid Inky's pincer, and
        deliberately uses teleporters to break a chase. On large
        \textsc{Tetris} mazes it routinely clears 80\% of the dots
        before being caught.
\end{itemize}

\subsection{Bot vs.\ ghost: the \texttt{RaceScreen}}

We added a dedicated screen, \texttt{RaceScreen}, that runs $N$ games
back-to-back with a fixed (algorithm, size, ghost difficulty, bot
strategy) tuple. Each game ends on Pac-Man's death or when the dots
are cleared, and the screen logs:
\begin{itemize}[leftmargin=1.3em,itemsep=1pt]
  \item average dots eaten per game,
  \item average game duration,
  \item bot decision time (BFS / Dijkstra),
  \item average A* nodes evaluated per ghost decision.
\end{itemize}
This is what produced the numbers in the tables of \S3 and \S4. Without
this offline harness, comparing strategies would have collapsed into
``it feels harder''.

% ---------- 5 ----------
\section{Rendering, UI and Networking}

These layers do not contain AI \emph{logic}, but they shape what the AI
\emph{looks like} and \emph{how it is debugged}.

\subsection{Rendering: \texttt{MazeRenderer}}

The renderer uses LibGDX's \texttt{ShapeRenderer} with a Pac-Man
inspired neon palette (deep blue walls, warm yellow dots, magenta
teleporters). Three details are worth highlighting from the AI
perspective:

\begin{itemize}[leftmargin=1.3em,itemsep=1pt]
  \item \textbf{Visual interpolation.} Both ghost and Pac-Man positions
        are rendered from \texttt{visualX/visualY} fields that
        \texttt{lerp} toward the logical \texttt{targetX/targetY} every
        frame. The player perceives a smooth 60\,FPS game even though
        the AI itself ticks at $\approx 5$\,Hz.
  \item \textbf{Visited trail.} A \texttt{boolean[][]\,visitedCells}
        grid is rendered as a slightly lighter background. This is a
        gameplay aid for human players, but it also doubles as a debug
        view for the bot: when the bot revisits cells too often, the
        trail makes it visually obvious.
  \item \textbf{Teleport links.} A faint translucent line is drawn
        between paired TP cells. This is purely cosmetic, but it helped
        us spot the oscillation bug of \S4.4: we could literally
        \emph{see} the bot ping-pong between two glowing pads.
\end{itemize}

\subsection{UI and screens}

The UI is structured as a small Scene2D state machine:
\texttt{MenuScreen} $\to$ \texttt{MazeScreen} (gameplay) $\to$
\texttt{ReplayListScreen} $\to$ \texttt{ReplayScreen}, with
\texttt{RaceScreen} as a separate benchmarking entry. Custom Neon
widgets (\texttt{NeonButton}, \texttt{NeonSlider}, \texttt{NeonToggle},
\texttt{NeonSkin}) provide a consistent visual language for the menu
and HUD. The HUD itself (\texttt{MazeHud}) overlays live AI telemetry
(ghost mode, A* nodes evaluated, FPS, score, lives) which made
balancing decisions data-driven instead of subjective.

\subsection{Networking: \texttt{ApiClient}}

\texttt{ApiClient} consumes the Flask backend asynchronously through
LibGDX's \texttt{Net} module. Two endpoints matter for AI behaviour:

\begin{itemize}[leftmargin=1.3em,itemsep=1pt]
  \item \texttt{GET /labyrinth} returns a freshly generated maze.
        Switching algorithm, height, width or \texttt{steps=true}
        produces different worlds for the same agents to play -- this
        is what makes the AI evaluation \emph{generalising} rather than
        memorising a fixed map.
  \item \texttt{POST /labyrinth/rate} feeds player feedback (0--5) back
        to MongoDB. Although we did not have time to use these ratings
        in a learning loop, the data is collected and persists across
        sessions, ready for offline analysis.
\end{itemize}

The client supports both a local \texttt{localhost:5000} backend and
the deployed \texttt{onrender.com} instance. The fallback is silent: if
the deployed instance times out, the game still starts (with a slightly
degraded UX -- no rating, no replay upload).

\subsection{Replay system}

\texttt{GameRecorder} captures every input frame -- direction key
presses for human runs, \texttt{(dr, dc)} pairs for the bot --
and \texttt{GameReplayer} feeds them back through a
\texttt{ReplayInputAgent} (an \texttt{Agent} implementation that
returns the recorded move). Combined with the deterministic seeding of
the maze (the maze ID identifies the exact grid), this gives us
\textbf{deterministic, repeatable AI runs}. Without it the
performance numbers in \S3 and \S4 would have been noise.

% ---------- 6 ----------
\section{AI Tools in Our Workflow}

Beyond the AI \emph{inside} the game, the supervisor's email explicitly
asked us to discuss \emph{how we used AI} during development. We took an
explicit, disciplined approach.

\subsection{What we used and how}

\begin{itemize}[leftmargin=1.3em,itemsep=1pt]
  \item \textbf{GitHub Copilot} (in-IDE completion) for boilerplate:
        getter/setter, JavaDoc skeletons, JSON parsing, repetitive
        Swagger annotations on Flask routes, LibGDX colour declarations.
  \item \textbf{Copilot Chat / ChatGPT} for higher-level scaffolding:
        first draft of A* pseudo-code, regex sanity checks, debugging
        hints when a stack trace was opaque, brainstorming the shape of
        the danger field.
  \item \textbf{Custom skill files} (\texttt{.github/skills/full-project-explorer})
        so the agent had a precise project map before suggesting edits.
        This single change drastically reduced wrong-context completions
        because the agent stopped guessing at file locations.
\end{itemize}

\subsection{Rules we set ourselves}

\begin{enumerate}[leftmargin=1.3em,itemsep=1pt]
  \item \textbf{Never accept code we cannot explain.} Every AI
        suggestion had to be read, understood, and rewritten in the
        project's style before being committed. This is the rule that
        caught the largest number of subtle bugs.
  \item \textbf{AI does not own the architecture.} The class hierarchy
        (\texttt{Ghost}, \texttt{Agent}, \texttt{BotAgent}, screen
        system), the four-algorithm split on the backend, and the API
        contract were designed by us. The AI was a \emph{typing
        accelerator}, not a \emph{decision maker}.
  \item \textbf{Tests stay human.} Critical correctness tests
        (\texttt{tests/test\_generators.py},
        \texttt{test\_maze\_properties.py}, the end-to-end shell tests)
        were written by hand to avoid the AI grading its own output.
  \item \textbf{Prompt with evidence.} When asking for help we always
        attached the actual file, never paraphrased it -- this avoided
        the classic ``hallucinated API'' problem and made the
        suggestions concrete.
  \item \textbf{Commit AI-assisted code separately.} Whenever a
        suggestion was non-trivial we committed it on its own so we
        could revert quickly if a regression appeared.
\end{enumerate}

\subsection{What worked, what did not}

\begin{table}[h]
\centering\small
\renewcommand{\arraystretch}{1.15}
\begin{tabularx}{\linewidth}{l X X}
\toprule
\textbf{Task} & \textbf{AI was helpful} & \textbf{AI was misleading} \\
\midrule
Flask + Swagger boilerplate & Saved hours of typing            & -- \\
A* on the maze grid         & Skeleton was correct             & Wrong wall indexing twice. \\
Ghost targeting maths       & Confirmed our formulas           & Suggested non-canonical Inky logic. \\
LibGDX rendering            & Helpful for \texttt{ShapeRenderer} usage & Outdated API hints from old major versions. \\
Teleport cooldown bug       & --                               & Suggested fixing it in the renderer (wrong layer). \\
Danger-field weights        & Suggested reasonable defaults    & Underestimated the impact of frightened ghosts. \\
\bottomrule
\end{tabularx}
\end{table}

\subsection{Take-away}

AI tools accelerated us most on \emph{shallow but tedious} work
(REST glue, JSON parsing, doc-blocks) and on \emph{recall} of well-known
algorithms (A*, BFS, Dijkstra). They were \textbf{least useful} on the
parts that defined the project: classical Pac-Man ghost personalities,
the maze symmetry pipeline, the danger-field calibration and the
timing-sensitive game logic. The lesson we take away is that AI
assistance scales with how well you can already articulate the problem
-- a good prompt is mostly a good specification.

% ---------- 7 ----------
\section{Conclusion}

The project demonstrates a layered approach to AI:

\begin{itemize}[leftmargin=1.3em,itemsep=1pt]
  \item A \textbf{procedurally generated world} (four maze algorithms
        with cycles, symmetry and teleports) that prevents any agent
        from memorising a fixed map.
  \item Four \textbf{classical ghost personalities} reimplemented from
        the original arcade specification, with both a faithful greedy
        mode and an instrumented A* mode for academic comparison.
  \item A \textbf{Pac-Man bot} that combines BFS dot-seeking with a
        Dijkstra danger field, including non-trivial details such as
        teleport cooldown and frightened-ghost masking.
  \item A \textbf{measurement harness} (RaceScreen + replay) that turns
        gameplay into reproducible, quantifiable runs.
\end{itemize}

Each layer feeds the next: the renderer hides the discrete tick of the
AI; the API unifies generation and rating; the replay system makes
evaluation reproducible. AI \emph{tools} accelerated the surrounding
scaffolding, but the design decisions that make the game play well --
targeting heuristics, teleport cooldown, danger-field weighting,
ghost-room exit logic -- remained firmly in our hands.

% ---------- References ----------
\begin{thebibliography}{9}
\small
\bibitem{pacmandossier} J.~Pittman. \emph{The Pac-Man Dossier}.
        \url{https://pacman.holenet.info/}, accessed 2026.
\bibitem{millington} I.~Millington, J.~Funge.
        \emph{Artificial Intelligence for Games}, 2nd ed.,
        Morgan Kaufmann, 2009.
\bibitem{russellnorvig} S.~Russell, P.~Norvig.
        \emph{Artificial Intelligence: A Modern Approach}, 4th ed.,
        Pearson, 2020.
\bibitem{libgdx} LibGDX Project. \url{https://libgdx.com/}.
\bibitem{flask} Pallets Projects. \emph{Flask documentation}.
        \url{https://flask.palletsprojects.com/}.
\end{thebibliography}

\end{document}
